\documentclass[9pt,landscape,a4paper]{extarticle}
\usepackage[utf8]{inputenc}
\usepackage[vietnamese]{babel}
\usepackage[T1]{fontenc}
\usepackage{fontspec}
\setmainfont{Libertinus Serif} % As per user's font

% Layout and multicolumn
\usepackage[margin=0.8cm, top=1.5cm, bottom=1.5cm]{geometry}
\usepackage{multicol}
\setlength{\columnsep}{0.8cm}
\setlength{\columnseprule}{0.1pt}

% Math and symbols
\usepackage{amsmath,amssymb}
\usepackage{wasysym}
\usepackage{xcolor}

% Custom Colors (Stanford Cheatsheet style)
\definecolor{stanfordblue}{RGB}{1, 104, 149}
\definecolor{stanfordred}{RGB}{140, 21, 21}
\definecolor{stanfordgreen}{RGB}{0, 100, 0}
\definecolor{cyanbox}{RGB}{220, 245, 255}

% Section styling (mimicking the sheet)
\usepackage{titlesec}
\titleformat{\section}
  {\normalfont\normalsize\bfseries}{\thesection}{1em}{}
\titlespacing*{\section}{0pt}{1.5ex plus 1ex minus .2ex}{0.5ex plus .2ex}

\titleformat{\subsection}
  {\normalfont\small\bfseries}{\thesubsection}{1em}{}
\titlespacing*{\subsection}{0pt}{1ex plus 1ex minus .2ex}{0.5ex plus .2ex}

% Fancy header
\usepackage{fancyhdr}
\pagestyle{fancy}
\fancyhf{}
\renewcommand{\headrulewidth}{0.4pt}
\renewcommand{\footrulewidth}{0.4pt}

% Header content
\lhead{\textsc{MAT3562E} -- \textsc{COMPUTER VISION}}
\rhead{\colorbox{cyan!20}{\makebox[4cm][c]{\textbf{Bùi Quang Chiến -- 23001837}}}}

% Footer content
\lfoot{\textsc{HUS -- VNU}}
\cfoot{\thepage}
\rfoot{\textsc{SPRING 2026}}

% Dense lists
\usepackage{enumitem}
\setlist{nosep, leftmargin=1.2em}
\setlist[itemize,1]{label={\textcolor{stanfordgreen}{\footnotesize$\square$}}}
\setlist[itemize,2]{label=--}

\begin{document}
\begin{center}
    {\Large \bfseries VIP Cheatsheet: LAB06} \\
    \vspace{0.3cm}
    {\large \textsc{Bùi Quang Chiến}} \\
    \vspace{0.2cm}
    {\small April 2026}
\end{center}
\vspace{0.3cm}

\begin{multicols*}{3}

\begin{center}
    {\Large \bfseries VIP Cheatsheet: LAB06} \\
    \vspace{0.3cm}
    {\large \textsc{Bùi Quang Chiến}} \\
    \vspace{0.2cm}
    {\small April 2026}
\end{center}
\vspace{0.3cm}

\begin{multicols*}{3}

\begin{center}
  {\large\bfseries\color{hdr} LAB 06 CHEATSHEET --- Region-Based Segmentation} \\
  {\small Bui Quang Chien \textbullet{} 23001837 \textbullet{} MAT3562E Computer Vision 2025--2026}
\end{center}
\hrule\vspace{2pt}

\begin{multicols}{3}

% ══════════════════════════════════════════════════════════════════════════════
\section{K-means --- Theory}

\textbf{Objective:} Partition $N$ pixels into $k$ clusters minimizing:
\[
  J = \sum_{j=1}^{k}\sum_{\mathbf{x}\in C_j}\lVert\mathbf{x}-\boldsymbol{\mu}_j\rVert^2
\]

\textbf{Algorithm (Lloyd's):}
\begin{enumerate}
  \item Init $k$ centroids (random or K-means++)
  \item \textbf{Assign:} $\text{label}(\mathbf{x})=\arg\min_j\lVert\mathbf{x}-\boldsymbol{\mu}_j\rVert^2$
  \item \textbf{Update:} $\boldsymbol{\mu}_j = \frac{1}{|C_j|}\sum_{\mathbf{x}\in C_j}\mathbf{x}$
  \item Repeat 2--3 until $\lVert\Delta\boldsymbol{\mu}\rVert < \varepsilon$
\end{enumerate}

\begin{notebox}
\textbf{K-means does NOT guarantee connected regions!}
Pixels with similar intensity far apart may share a label.
\end{notebox}

\begin{infobox}
\textbf{Feature vector variants:}\\
\textbf{Intensity only:} $\mathbf{x}=(I)$ \quad 1-D\\
\textbf{Intensity+XY:} $\mathbf{x}=(I,\; w\cdot y/H,\; w\cdot x/W)$ \quad 3-D\\
$w$ = \texttt{spatial\_weight} controls locality.
\end{infobox}

% ══════════════════════════════════════════════════════════════════════════════
\section{K-means --- Intensity Only}

\begin{lstlisting}[language=Python]
def kmeans_gray_intensity(gray, k=3,
      max_iter=30, tol=1e-4, seed=0):
  X = gray.astype(np.float32).reshape(-1, 1)
  rng = np.random.default_rng(seed)
  idx = rng.choice(X.shape[0], size=k,
                   replace=False)
  centers = X[idx].copy()

  for _ in range(max_iter):
    dists = np.sum(
      (X[:,None,:] - centers[None,:,:]) **2,
      axis=2)
    labels = np.argmin(dists, axis=1)

    new_c = centers.copy()
    for c in range(k):
      pts = X[labels == c]
      if len(pts) > 0:
        new_c[c] = pts.mean(axis=0)

    if np.linalg.norm(new_c - centers) < tol:
      break
    centers = new_c

  return labels.reshape(gray.shape),\
         centers.reshape(-1)
\end{lstlisting}

\begin{tipbox}
\textbf{Key idea:} Flatten image to $N\times 1$, cluster by intensity, reshape labels back.
\end{tipbox}

% ══════════════════════════════════════════════════════════════════════════════
\section{K-means --- Intensity + XY}

\begin{lstlisting}[language=Python]
def kmeans_gray_intensity_xy(gray, k=3,
      spatial_weight=0.25, max_iter=30,
      tol=1e-4, seed=0):
  h, w = gray.shape
  yy, xx = np.mgrid[0:h, 0:w]
  I = gray.astype(np.float32) / 255.0
  X = np.stack([
    I,
    spatial_weight * yy.astype(np.float32)
                     / max(h-1, 1),
    spatial_weight * xx.astype(np.float32)
                     / max(w-1, 1)
  ], axis=-1).reshape(-1, 3)
  # ... same loop as intensity-only ...
\end{lstlisting}

\begin{center}
\begin{tabular}{@{}cl@{}}
\toprule
\texttt{spatial\_weight} & \textbf{Effect} \\
\midrule
0.0 & Pure intensity (= basic K-means)\\
0.25 & Slight spatial coherence\\
0.5--0.7 & Strong locality, like superpixels\\
$>1.0$ & Ignores intensity, grid-like\\
\bottomrule
\end{tabular}
\end{center}

% ══════════════════════════════════════════════════════════════════════════════
\section{K-means --- OpenCV}

\begin{lstlisting}[language=Python]
def kmeans_opencv_intensity(gray, k=2,
                            attempts=10):
  X = gray.reshape(-1,1).astype(np.float32)
  criteria = (cv2.TERM_CRITERIA_EPS +
              cv2.TERM_CRITERIA_MAX_ITER,
              50, 0.2)
  comp, labels, centers = cv2.kmeans(
      X, k, None, criteria, attempts,
      cv2.KMEANS_PP_CENTERS)
  return labels.reshape(gray.shape),\
         centers.reshape(-1), comp
\end{lstlisting}

\begin{tipbox}
\textbf{Advantages:} C++ speed, K-means++ init, multiple \texttt{attempts} $\to$ more stable.
\texttt{compactness} = total within-cluster distance (lower is better).
\end{tipbox}

% ══════════════════════════════════════════════════════════════════════════════
\section{Region Growing --- Theory}

\textbf{Idea:} BFS/DFS from seed pixel; add neighbor if it meets homogeneity criterion.

\textbf{Two reference strategies:}

\begin{infobox}
\textbf{R1 --- Seed Reference:}\\
$|I(x,y) - I_{\text{seed}}| \leq T$\\
Stable but misses gradual gradients.
\end{infobox}

\begin{infobox}
\textbf{R2 --- Region Mean:}\\
$|I(x,y) - \bar{I}_{\text{region}}| \leq T$\\
Adapts to gradient but may ``leak''.
\end{infobox}

\textbf{Connectivity:}
\begin{itemize}
  \item \textbf{4-conn:} up/down/left/right
  \item \textbf{8-conn:} + 4 diagonals (more natural but leaks easier)
\end{itemize}

% ══════════════════════════════════════════════════════════════════════════════
\section{Region Growing --- Seed Ref}

\begin{lstlisting}[language=Python]
from collections import deque

def region_growing_seed_ref(gray, seed,
      threshold=10, connectivity=4):
  h, w = gray.shape
  visited = np.zeros((h,w), dtype=bool)
  region  = np.zeros((h,w), dtype=np.uint8)
  ref_val = float(gray[seed])

  nbrs = [(-1,0),(1,0),(0,-1),(0,1)]
  if connectivity == 8:
    nbrs += [(-1,-1),(-1,1),(1,-1),(1,1)]

  q = deque([seed]); visited[seed] = True
  while q:
    x, y = q.popleft()
    if abs(float(gray[x,y]) - ref_val)\
       <= threshold:
      region[x,y] = 255
      for dx,dy in nbrs:
        nx, ny = x+dx, y+dy
        if 0<=nx<h and 0<=ny<w\
           and not visited[nx,ny]:
          visited[nx,ny] = True
          q.append((nx,ny))
  return region
\end{lstlisting}

% ══════════════════════════════════════════════════════════════════════════════
\section{Region Growing --- Region Mean}

\begin{lstlisting}[language=Python]
def region_growing_region_mean(gray, seed,
      threshold=10, connectivity=4):
  h, w = gray.shape
  visited = np.zeros((h,w), dtype=bool)
  region  = np.zeros((h,w), dtype=np.uint8)
  region_sum = float(gray[seed])
  region_count = 1
  # ... same BFS but:
  # ref_val = region_sum / region_count
  # when adding pixel:
  #   region_sum += float(gray[nx,ny])
  #   region_count += 1
\end{lstlisting}

\begin{notebox}
\textbf{Leak risk:} Region mean drifts as the region grows. Use with smaller $T$ near boundaries.
\end{notebox}

% ══════════════════════════════════════════════════════════════════════════════
\section{Key Parameters}

\begin{center}
\begin{tabular}{@{}lll@{}}
\toprule
\textbf{Param} & \textbf{Method} & \textbf{Effect} \\
\midrule
$k$ & K-means & \# of segments \\
& & $\uparrow k$: more detail\\
\texttt{w} & K-means & spatial weight\\
& +XY & $\uparrow w$: more local\\
$T$ & Region & tolerance\\
& Growing & $\uparrow T$: larger region\\
conn & Region & 4 vs 8\\
& Growing & 8: smoother, leakier\\
seed & Region & start point\\
& Growing & determines which region\\
\bottomrule
\end{tabular}
\end{center}

% ══════════════════════════════════════════════════════════════════════════════
\section{Common Mistakes}

\begin{notebox}
\begin{enumerate}
  \item \textbf{No preprocessing:} Noise makes K-means and Region Growing unreliable. Always consider Gaussian blur.
  \item \textbf{Wrong dtype:} OpenCV \texttt{kmeans} requires \texttt{float32}.
  \item \textbf{Overflow in distance:} Use \texttt{float32} before squaring.
  \item \textbf{Forgetting visited:} Region Growing without visited array $\to$ infinite loop.
  \item \textbf{Seed outside object:} Region Growing result depends entirely on seed position.
\end{enumerate}
\end{notebox}

\end{multicols}

\end{multicols*}

\end{multicols*}
\end{document}
